\documentclass[11pt]{article}
\usepackage[T1]{fontenc}
\usepackage{amsmath,amssymb,amsthm,lmodern}
\usepackage[margin=1in]{geometry}
\usepackage{hyperref}
\hypersetup{hidelinks,hypertexnames=false,pdftitle={MCMT: Broder Card-Marking Strong Stationary Time, Proposition 8.6},pdfauthor={David A. Levin; Yuval Peres}}
\newcommand{\X}{\mathcal{X}}
\newcommand{\Ft}{\mathcal{F}_t}
\newcommand{\Px}{\mathbf{P}_x}
\newcommand{\sourceloc}[1]{\par\smallskip\noindent\textit{Source: #1.}\par\nopagebreak[4]\smallskip}
\newcommand{\editorial}[1]{\par\smallskip\noindent\textsf{Editorial note.} #1\par\smallskip}
\renewcommand{\qedsymbol}{\ensuremath{\blacksquare}}
\title{Broder Card-Marking Strong Stationary Time\\Proposition 8.6}
\author{David A. Levin \and Yuval Peres}
\date{Private reading transcription of the fixed second-edition author PDF}
\begin{document}
\maketitle

% BEGIN CONTEXT MCMT:8.6
\noindent\textbf{Source and selected scope.}
\emph{Markov Chains and Mixing Times, second edition}, by David A. Levin
and Yuval Peres, ``With contributions by Elizabeth L. Wilmer'' as printed
on physical page 1. The fixed author-hosted PDF has 461 physical pages
and PDF generation metadata dated March 21, 2017; a separate author
revision date is not stated on the inspected cover.
The source is \url{https://pages.uoregon.edu/dlevin/MARKOV/mcmt2e.pdf}.
The sole primary is Proposition 8.6 and its complete printed proof on
physical pages 121--122, printed pages 105--106. The stacking paragraph
is retained as the author's heuristic. Definitions, the transition
rule, Remark 8.7, complete Lemma 8.9 and Corollary 8.10, and the selected
Notes are uncounted context or consequences. This transcription remains
private.

\subsection*{Selected stopping-time definitions and mixing convention}
\sourceloc{Section 6.2, physical page 93, printed page 77}
Let $\{\Ft\}$ be a \emph{filtration}, a sequence of $\sigma$-algebras such
that $\Ft\subset\mathcal{F}_{t+1}$ for all $t$. We say that $\{X_t\}$ is
\emph{adapted} to $\{\Ft\}$ if $X_t$ is $\Ft$-measurable for all $t$.
If $\mathcal{H}_t=\sigma(X_0,X_1,\ldots,X_t)$, then
$\{\mathcal{H}_t\}$ is called the \emph{natural filtration}. Clearly
$\{X_t\}$ is adapted to the natural filtration.

Suppose that $\{X_t\}$ is adapted to $\{\Ft\}$. We say that $\{X_t\}$ is
a \emph{Markov chain with respect to $\{\Ft\}$} if
\begin{equation}
\Px\{X_{t+1}=y\mid\Ft\}=P(X_t,y),
\tag{6.1}\label{eq:6.1}
\end{equation}
where $P$ is a transition matrix. A Markov chain as defined by (1.1)
satisfies \eqref{eq:6.1} when $\{\Ft\}$ is the natural filtration.

A \emph{stopping time} for the filtration $\{\Ft\}$ is a random variable
$\tau$ with values in $\{0,1,\ldots\}$ such that $\{\tau=t\}\in\Ft$.
If the filtration of a stopping time is not specified, it will be assumed
to be the natural filtration.

\sourceloc{Section 6.4, physical page 94, printed page 78; strong stationary time definition only}
Let $(X_t)$ be a Markov chain with respect to the filtration $\{\Ft\}$,
with stationary distribution $\pi$. A \emph{strong stationary time} for
$(X_t)$ and starting position $x$ is an $\{\Ft\}$-stopping time $\tau$,
such that for all times $t$ and all $y$,
\begin{equation}
\Px\{\tau=t,X_\tau=y\}=\Px\{\tau=t\}\pi(y).
\tag{6.3}\label{eq:6.3}
\end{equation}
In words, $X_\tau$ has distribution $\pi$ and is independent of $\tau$.


\sourceloc{Proposition 6.11, physical page 95, printed page 79; statement used as an uncounted input}
\noindent\textbf{Proposition 6.11.}
\emph{If $\tau$ is a strong stationary time for starting state $x$, then}
\begin{equation}
\|P^t(x,\cdot)-\pi\|_{\mathrm{TV}}\le\Px\{\tau>t\}.
\tag{6.5}\label{eq:6.5}
\end{equation}

\sourceloc{Selected conventions from Sections 4.4--4.5, physical pages 69--70, printed pages 53--54}
It is therefore convenient to define
\begin{equation}
d(t):=\max_{x\in\X}\|P^t(x,\cdot)-\pi\|_{\mathrm{TV}}.
\tag{4.22}\label{eq:4.22}
\end{equation}
The mixing time is defined by
\begin{equation}
t_{\mathrm{mix}}(\varepsilon):=\min\{t:d(t)\le\varepsilon\}
\tag{4.30}\label{eq:4.30}
\end{equation}
and
\begin{equation}
t_{\mathrm{mix}}:=t_{\mathrm{mix}}(1/4).
\tag{4.31}\label{eq:4.31}
\end{equation}
% END CONTEXT MCMT:8.6

% BEGIN SHUFFLE CONTEXT MCMT:8.6
\subsection*{The random transposition shuffle}
\sourceloc{Section 8.1, physical page 116, printed page 100; selected convention}
A probability distribution $\mu$ on the symmetric group describes a
mechanism for shuffling cards: apply permutation $\sigma$ to the deck
with probability $\mu(\sigma)$. Repeatedly shuffling the deck using this
mechanism is equivalent to running the random walk on the group with
increment distribution $\mu$. As discussed in Section 2.6, as long as
the support of $\mu$ generates all of $S_n$, the resulting chain is
irreducible. If $\mu(\mathrm{id})>0$, then it is aperiodic. Every shuffle
chain has uniform stationary distribution.

\sourceloc{Section 8.2, physical page 118, printed page 102; complete transition rule and equation (8.3)}
To avoid periodicity, the random transposition shuffle is defined as
follows: At time $t$, choose two cards, labelled $L_t$ and $R_t$,
independently and uniformly at random. If $L_t$ and $R_t$ are different,
transpose them. Otherwise, do nothing. The resulting distribution $\mu$
satisfies
\begin{equation}
\mu(\sigma)=\begin{cases}
1/n&\text{if }\sigma=\mathrm{id},\\
2/n^2&\text{if }\sigma=(ij),\\
0&\text{otherwise}.
\end{cases}
\tag{8.3}\label{eq:8.3}
\end{equation}
As in Example 8.2, this chain is irreducible; aperiodicity follows from
Remark 8.3.

\editorial{The earlier Section 8.2 lower-bound discussion says,
``Consider the random transposition shuffle started from $\mathrm{id}$.''
The primary proposition below is transcribed with its own printed
initialization: no cards have been marked at time $0$. No additional
claim about an arbitrary random initial law is inserted.}
% END SHUFFLE CONTEXT MCMT:8.6

\section*{8.2.3. Upper bound via strong stationary time}
\sourceloc{Physical page 121, printed page 105; complete Proposition 8.6}
% BEGIN PRIMARY STATEMENT MCMT:8.6
\noindent\textbf{Proposition 8.6.}
\emph{In the random transposition shuffle, let $R_t$ and $L_t$ be the
cards chosen by the right and left hands, respectively, at time $t$.
Assume that when $t=0$, no cards have been marked. At time $t$, mark card
$R_t$ if both of the following are true:}
\begin{itemize}
\item $R_t$ is unmarked.
\item Either $L_t$ is a marked card or $L_t=R_t$.
\end{itemize}
\emph{Let $\tau$ be the time when every card has been marked. Then $\tau$
is a strong stationary time for this chain.}
% END PRIMARY STATEMENT MCMT:8.6

% BEGIN HEURISTIC CONTEXT MCMT:8.6
Here is a heuristic explanation for why the scheme described above
should give a strong stationary time. One way to generate a uniform
random permutation is to build a stack of cards, one at a time,
inserting each card into a uniformly random position relative to the
cards already in the stack. For the stopping time described above, the
marked cards are carrying out such a process.
% END HEURISTIC CONTEXT MCMT:8.6

\sourceloc{Physical pages 121--122, printed pages 105--106; complete printed proof, stopping before Remark 8.7}
% BEGIN PRIMARY PROOF MCMT:8.6
\begin{proof}
It is clear that $\tau$ is a stopping time. To show that it is a strong
stationary time, we prove the following subclaim by induction on $t$.
Let $V_t\subseteq[n]$ be the set of cards marked at or before time $t$,
and let $U_t\subseteq[n]$ be the set of positions occupied by $V_t$
after the $t$-th transposition. We claim that given $t$, $V_t$, and
$U_t$, all possible permutations of the cards in $V_t$ on the positions
$U_t$ are equally likely.

This is clearly true when $t=1$ (and continues to be true as long as at
most one card has been marked).

Now, assume that the subclaim is true for $t$. The shuffler chooses
cards $L_{t+1}$ and $R_{t+1}$.
\begin{itemize}
\item If no new card is marked, then $V_{t+1}=V_t$. This can happen in
three ways:
\begin{itemize}
\item The cards $L_{t+1}$ and $R_{t+1}$ are different and both are
unmarked. Then $V_{t+1}$ and $U_{t+1}$ are identical to $V_t$ and $U_t$,
respectively.
\item If $L_{t+1}$ and $R_{t+1}$ were both marked at an earlier round,
then $U_{t+1}=U_t$ and the shuffler applies a uniform random
transposition to the cards in $V_t$. All permutations of $V_t$ remain
equiprobable.
\item Otherwise, $L_{t+1}$ is unmarked and $R_{t+1}$ was marked at an
earlier round. To obtain the position set $U_{t+1}$, we delete the
position (at time $t$) of $R_{t+1}$ and add the position (at time $t$)
of $L_{t+1}$. For a fixed set $U_t$, all choices of $R_{t+1}\in U_t$
are equally likely, as are all permutations of $V_t$ on $U_t$. Hence,
once the positions added and deleted are specified, all permutations
of $V_t$ on $U_{t+1}$ are equally likely.
\end{itemize}
\item If the card $R_{t+1}$ gets marked, then $L_{t+1}$ is equally
likely to be any element of $V_{t+1}=V_t\cup\{R_{t+1}\}$, while
$U_{t+1}$ consists of $U_t$ along with the position of $R_{t+1}$
(at time $t$). Specifying the permutation of $V_t$ on $U_t$ and the
card $L_{t+1}$ uniquely determines the permutation of $V_{t+1}$ on
$U_{t+1}$. Hence all such permutations are equally likely.
\end{itemize}
In every case, the collection of all permutations of the cards $V_t$
on a specified set $U_t$ together make equal contributions to all
possible permutations of $V_{t+1}$ on $U_{t+1}$. Hence, to conclude that
all possible permutations of a fixed $V_{t+1}$ on a fixed $U_{t+1}$
are equally likely, we simply sum over all possible preceding
configurations.
\end{proof}
% END PRIMARY PROOF MCMT:8.6

\sourceloc{Physical page 122, printed page 106; complete Remark 8.7}
% BEGIN REMARK CONTEXT MCMT:8.6
\noindent\textbf{Remark 8.7.}
In the preceding proof, the two subcases of the inductive step for
which no new card is marked are essentially the same as checking that
the uniform distribution is stationary for the random transposition
shuffle and the random-to-top shuffle, respectively.

\editorial{The primary proof preserves the printed phrase
$R_{t+1}\in U_t$, while the preceding definition calls $V_t$ the card
set and $U_t$ the position set. This notation is retained literally.
The original base-case and predecessor-summation abbreviations have
not been expanded.}
% END REMARK CONTEXT MCMT:8.6

\subsection*{Complete moment and mixing consequences, uncounted}
\sourceloc{Physical pages 122--123, printed pages 106--107; complete Lemma 8.9 and its printed proof}
% BEGIN SUPPORT STATEMENT MCMT:8.9
\noindent\textbf{Lemma 8.9.}
\emph{The stopping time $\tau$ defined in Proposition 8.6 satisfies}
\[
\mathbf{E}(\tau)=2n(\log n+O(1))
\]
\emph{and}
\[
\operatorname{Var}(\tau)=O(n^2).
\]
% END SUPPORT STATEMENT MCMT:8.9
% BEGIN SUPPORT PROOF MCMT:8.9
\begin{proof}
As for the coupon collector time, we can decompose
\[
\tau=\tau_0+\cdots+\tau_{n-1},
\]
where $\tau_k$ is the number of transpositions after the $k$-th card
is marked, up to and including when the $(k+1)$-st card is marked.
The rules specified in Proposition 8.6 imply that $\tau_k$ is a
geometric random variable with success probability
$\frac{(k+1)(n-k)}{n^2}$ and that the $\tau_i$'s are independent of
each other. Hence
\[
\mathbf{E}(\tau)=\sum_{k=0}^{n-1}\frac{n^2}{(k+1)(n-k)}.
\]
Substituting the partial fraction decomposition
\[
\frac{1}{(k+1)(n-k)}
=\frac{1}{n+1}\left[\frac{1}{k+1}+\frac{1}{n-k}\right]
\]
and recalling that
\[
\sum_{j=1}^n\frac1j=\log n+O(1)
\]
(see Exercise 2.4) completes the estimate.

Now, for the variance. We can immediately write
\[
\operatorname{Var}(\tau)
=\sum_{k=0}^{n-1}
\frac{1-\frac{(k+1)(n-k)}{n^2}}
{\left(\frac{(k+1)(n-k)}{n^2}\right)^2}
<\sum_{k=0}^{n-1}\frac{n^4}{(k+1)^2(n-k)^2}.
\]
Split the sum into two pieces:
\begin{align*}
\operatorname{Var}(\tau)
&<\sum_{0\le k<n/2}\frac{n^4}{(k+1)^2(n-k)^2}
 +\sum_{n/2\le k<n}\frac{n^4}{(k+1)^2(n-k)^2}\\
&<\frac{2n^4}{(n/2)^2}
\sum_{0\le k\le n/2}\frac{1}{(k+1)^2}=O(n^2).
\end{align*}
\end{proof}
% END SUPPORT PROOF MCMT:8.9

\sourceloc{Physical page 123, printed page 107; complete Corollary 8.10 and its printed proof}
% BEGIN SUPPORT STATEMENT MCMT:8.10
\noindent\textbf{Corollary 8.10.}
\emph{For the random transposition chain on an $n$-card deck,}
\[
t_{\mathrm{mix}}\le(2+o(1))n\log n.
\]
% END SUPPORT STATEMENT MCMT:8.10
% BEGIN SUPPORT PROOF MCMT:8.10
\begin{proof}
Let $\tau$ be the Broder stopping time defined in Proposition 8.6, and
let $t_0=\mathbf{E}(\tau)+2\sqrt{\operatorname{Var}(\tau)}$.
By Chebyshev's inequality,
\[
\mathbf{P}\{\tau>t_0\}\le\frac14.
\]
Lemma 8.9 and Proposition 6.11 now imply the desired inequality.
\end{proof}
% END SUPPORT PROOF MCMT:8.10

% BEGIN CREDIT CONTEXT MCMT:8.6
\sourceloc{Selected original Notes, physical page 128, printed page 112}
The strong stationary time defined in Proposition 8.6 and used to
prove the upper bound on the mixing time for random transpositions
(Corollary 8.10) is due to A. Broder (see Diaconis (1988a)). This upper
bound is off by a factor of 4. Matthews (1988b) gives an improved
strong stationary time whose upper bound matches the lower bound.

\editorial{These are the book's attribution and comparison statements.
The book authors provide the collected exposition; Broder's route
credit is kept separately. The later Matthews construction and its
unexpanded analysis are outside this selected primary. No historical
article or monograph has been newly acquired for this transcription.}
% END CREDIT CONTEXT MCMT:8.6
\end{document}
